% LaTeX port of resume.html, rendered through Jinja2 with custom tag delimiters.
\documentclass[10pt]{extarticle}

\usepackage[a4paper, left=19.05mm, right=19.05mm, top=14.53mm, bottom=32mm, footskip=29pt, headheight=16pt, headsep=10pt]{geometry}
\usepackage{fontspec}
\setmainfont{Poppins}[Path=./fonts/, Extension=.ttf, UprightFont=*-Regular, BoldFont=*-Bold]
\newfontfamily\popregular{Poppins}[Path=./fonts/, Extension=.ttf, UprightFont=*-Regular]
\newfontfamily\popmedium{Poppins}[Path=./fonts/, Extension=.ttf, UprightFont=*-Medium]
\newfontfamily\popsemibold{Poppins}[Path=./fonts/, Extension=.ttf, UprightFont=*-SemiBold]

\usepackage{xcolor}
\definecolor{inkblack}{HTML}{0F1828}
\definecolor{inkgray}{HTML}{57565D}
\definecolor{bgwhite}{HTML}{F2F2F2}

\usepackage{qrcode}
\usepackage{fontawesome5}
\usepackage{graphicx}
\usepackage{calc}
\usepackage{enumitem}
\usepackage{tabularx}
\usepackage{fancyhdr}
\usepackage{hyperref}
\hypersetup{colorlinks=false, pdfborder={0 0 0}, pdftitle={\VAR{ data.name|tex } - \VAR{ data.role|tex }}}
\usepackage{parskip}
\setlength{\parindent}{0pt}
\pagecolor{bgwhite}
\setlength{\parskip}{0pt}
\hyphenpenalty=10000
\exhyphenpenalty=10000

\setlist[itemize]{nosep, leftmargin=9pt, topsep=0pt, itemsep=1.5pt, label={\textbullet}}

% ---- Semantic text styles (sizes converted from the HTML template's px at 0.75pt/px) ----
\newcommand{\titleOne}[1]{{\popsemibold\fontsize{19.5}{29.25}\selectfont\color{inkblack}#1}}
\newcommand{\titleTwo}[1]{{\popregular\fontsize{16.5}{24.75}\selectfont\color{inkblack}#1}}
\newcommand{\titleThree}[1]{{\popsemibold\fontsize{11.25}{15.75}\selectfont\color{inkblack}#1}}
\newcommand{\labelText}[1]{{\popmedium\fontsize{9}{15.75}\selectfont\color{inkblack}#1}}
\newcommand{\paraText}[1]{{\popregular\fontsize{8.25}{13.5}\selectfont\color{inkgray}#1}}
\newcommand{\paraContentText}[1]{{\popregular\fontsize{8.25}{15}\selectfont\color{inkgray}#1}}
\newcommand{\paraBlack}[1]{{\popregular\fontsize{8.25}{13.5}\selectfont\color{inkblack}#1}}
\newcommand{\sectiondivider}{\par\vspace{2pt}{\color{inkblack}\rule{12.75pt}{1.5pt}}\vspace{6pt}}

% ---- Two-column layouts: section title|content, and entry date|content ----
\newlength{\sectiontitlewidth}\setlength{\sectiontitlewidth}{0.2415\textwidth}
\newlength{\sectioncontentwidth}\setlength{\sectioncontentwidth}{0.7431\textwidth}
\newlength{\entrydatewidth}\setlength{\entrydatewidth}{0.1657\sectioncontentwidth}
\newlength{\entrycontentwidth}\setlength{\entrycontentwidth}{0.8\sectioncontentwidth}

\newcommand{\resumesection}[2]{%
  \noindent\begin{minipage}[t]{\sectiontitlewidth}#1\end{minipage}\hfill%
  \begin{minipage}[t]{\sectioncontentwidth}\raggedright #2\end{minipage}\par\vspace{18pt}%
}
\newcommand{\entryrow}[2]{%
  \noindent\begin{minipage}[t]{\entrydatewidth}\paraContentText{#1}\end{minipage}\hfill%
  \begin{minipage}[t]{\entrycontentwidth}#2\end{minipage}\par\vspace{12pt}%
}
% A minipage can't break across a page, including mid-entry (one job with many
% bullets). resume.html's JS pagination broke between bullets to pack pages
% tightly; to match that here, Professional Experience doesn't box its entries
% at all -- it flows as ordinary indented paragraphs (\leftskip + a hanging
% date label), so LaTeX's own page breaker can split anywhere a paragraph or
% list normally could, same as every other multi-page LaTeX document.
\newcommand{\sectiontitleonly}[1]{\noindent\begin{minipage}[t]{\sectiontitlewidth}#1\end{minipage}\par\vspace{10pt}}
\newlength{\experienceindent}\setlength{\experienceindent}{\dimexpr\sectiontitlewidth+7.5pt+\entrydatewidth+7.5pt\relax}
\newcommand{\entrystart}[1]{\hspace*{-\experienceindent}\hspace{\sectiontitlewidth+7.5pt}%
  \begin{minipage}[t]{\entrydatewidth}\paraContentText{#1}\end{minipage}\hspace{7.5pt}}
\newenvironment{experienceentry}{\par\begingroup\raggedright\leftskip=\experienceindent\relax\parindent=0pt\setlength{\parskip}{0pt}\noindent}{\endgroup\vspace{12pt}}

% ---- Running headers/footers ----
\newcommand{\socialicons}{{\fontsize{14}{14}\selectfont\color{inkblack}%
  \href{\VAR{ data.contact.github }}{\faGithub}\hspace{16pt}\href{\VAR{ data.contact.linkedin }}{\faLinkedinIn}}}
\newcommand{\footerrow}{%
  \paraText{\VAR{ data.name|tex }\ \textcolor{inkblack}{Resume | P\thepage}}\hfill%
  \href{\VAR{ data.contact.website_link }}{\paraBlack{\VAR{ data.contact.website|tex }}}%
}
\fancypagestyle{firstpage}{\fancyhf{}\renewcommand{\headrulewidth}{0pt}\renewcommand{\footrulewidth}{0pt}%
  \fancyfoot[L,R]{}\fancyfoot[C]{}%
  \fancyfootoffset[L]{0pt}%
  \fancyfoot[L]{\footerrow}%
}
\fancypagestyle{restpage}{\fancyhf{}\renewcommand{\headrulewidth}{0pt}\renewcommand{\footrulewidth}{0pt}%
  \fancyhead[L]{\includegraphics[width=30pt,height=12pt]{page-logo.pdf}\hfill\socialicons}%
  \fancyfoot[L]{\footerrow}%
}
\pagestyle{restpage}

\begin{document}
\thispagestyle{firstpage}

% ---- Page 1 header: logo, name/role, social, contact, QR ----
% [t] aligns minipages by their first line's BASELINE, not their visual top; since
% the logo (12pt tall) and QR (72pt tall) differ a lot in height, that leaves their
% top edges misaligned unless the taller one (QR) is lowered by the difference.
\noindent\begin{minipage}[t]{\textwidth-78pt-18pt}
  \includegraphics[width=30pt,height=12pt]{header-logo.pdf}\par\vspace{18pt}
  \titleOne{\VAR{ data.name|tex }}\par
  \titleTwo{\VAR{ data.role|tex }}\par\vspace{10pt}
  \socialicons\par\vspace{10pt}
  \paraText{\href{\VAR{ data.contact.phone_link }}{\VAR{ data.contact.phone|tex }}}\par\vspace{2pt}
  \paraText{\href{mailto:\VAR{ data.contact.email }}{\VAR{ data.contact.email|tex }}}\par\vspace{2pt}
  \paraText{\VAR{ data.contact.location|tex }}\par\vspace{2pt}
  \paraText{\href{\VAR{ data.contact.website_link }}{\VAR{ data.contact.website|tex }}}
\end{minipage}\hfill
\raisebox{-60pt}{\begin{minipage}[t]{78pt}
  \centering\href{\VAR{ data.contact.qr_url }}{\qrcode[height=72pt, level=M]{\VAR{ data.contact.qr_url }}}
\end{minipage}}
\par\vspace{24pt}

\BLOCK{ if data.summary }
\resumesection{\titleThree{Summary\\of Experience}\sectiondivider}{\paraText{\VAR{ data.summary|tex }}}
\BLOCK{ endif }

\BLOCK{ if data.skills }
\resumesection{\titleThree{Technical\\Proficiencies}\sectiondivider}{%
  \BLOCK{ for group in data.skills }
  \begin{tabularx}{\sectioncontentwidth}{@{}XXXX@{}}
  \BLOCK{ for row in group|batch(4) }
  \BLOCK{ set cells = row + ['']*(4 - row|length) }
  \paraText{\VAR{ cells[0]|tex }} & \paraText{\VAR{ cells[1]|tex }} & \paraText{\VAR{ cells[2]|tex }} & \paraText{\VAR{ cells[3]|tex }} \\[4pt]
  \BLOCK{ endfor }
  \end{tabularx}\par\vspace{6pt}
  \BLOCK{ endfor }
}
\BLOCK{ endif }

\BLOCK{ if data.experience }
\sectiontitleonly{\titleThree{Professional\\Experience}\sectiondivider}
\BLOCK{ for job in data.experience }
\BLOCK{ set parts = experience_date_parts(job.dates) }
\begin{experienceentry}
\entrystart{\VAR{ parts[0]|tex }\BLOCK{ if parts[1] }\\--\ \VAR{ parts[1]|tex }\BLOCK{ endif }}%
  \labelText{\VAR{ job.company|tex }}\BLOCK{ if job.detail } \paraText{\VAR{ job.detail|tex }}\BLOCK{ endif }\par
  \BLOCK{ if job.role }\paraText{\VAR{ job.role|tex }}\par\BLOCK{ endif }
  \vspace{4pt}
  \BLOCK{ if job.description }\paraText{\VAR{ job.description|tex }}\par\vspace{2pt}\BLOCK{ endif }
  \BLOCK{ if job.bullets }
  \begin{itemize}[leftmargin=\dimexpr\experienceindent+9pt\relax, itemindent=0pt]
  \BLOCK{ for bullet in job.bullets }
  \item \paraText{\VAR{ bullet|tex }}
  \BLOCK{ endfor }
  \end{itemize}
  \BLOCK{ endif }
\end{experienceentry}
\BLOCK{ endfor }
\BLOCK{ endif }

\BLOCK{ if data.education }
\resumesection{\titleThree{Academic\\Education}\sectiondivider}{%
  \BLOCK{ for entry in data.education }
  \entryrow{\VAR{ entry.dates|tex }}{\labelText{\VAR{ entry.degree|tex }}\par\paraText{\VAR{ entry.institution|tex }}}
  \BLOCK{ endfor }
}
\BLOCK{ endif }

\BLOCK{ if data.languages }
\resumesection{\titleThree{Languages}\sectiondivider}{%
  \begin{tabularx}{\sectioncontentwidth}{@{}XXX@{}}
  \BLOCK{ for row in data.languages|batch(3) }
  \BLOCK{ set cells = row + [none, none, none] }
  \BLOCK{ for cell in cells[:3] }
  \BLOCK{ if cell }\paraText{\textcolor{inkblack}{\VAR{ cell.language|tex }}\ -\ \VAR{ cell.level|tex }}\BLOCK{ else }~\BLOCK{ endif }\BLOCK{ if not loop.last } & \BLOCK{ endif }
  \BLOCK{ endfor }
  \\[4pt]
  \BLOCK{ endfor }
  \end{tabularx}
}
\BLOCK{ endif }

\end{document}
